% !TEX root = arpa_e.tex
% =============================================================
%  Chapter — System Studies (HVDC design process)
%  Four study packages in design order:
%    steady state -> transient -> dynamic -> RMS
%  Figures: fig/design-process.tex (inline TikZ, \input below).
% =============================================================
\newcommand{\ssfig}[1]{fig/#1}

% \ssstudytable{tools}{time step}{main results} — key data of one study
\newcommand{\ssstudytable}[3]{%
  {\small\renewcommand{\arraystretch}{1.35}%
  \begin{tabularx}{\linewidth}{@{}l>{\raggedright\arraybackslash}X@{}}
    \toprule
    \atrowfill\textbf{Tools}     & #1 \\
    \textbf{Time step}           & #2 \\
    \atrowfill\textbf{Results}   & #3 \\
    \bottomrule
  \end{tabularx}}%
}

\sectionslide{System Studies}{HVDC design process and simulation packages}

\begin{frame}{Abbreviations}
  \begin{columns}[T,onlytextwidth]
    \begin{column}{0.48\textwidth}
      {\small
      \begin{tabularx}{\linewidth}{@{}l@{\hspace{0.8em}}X@{}}
        AC    & Alternating current \\
        DC    & Direct current \\
        EMT   & Electromagnetic transient \\
        EMTP  & Electromagnetic Transients Program \\
        FRT   & Fault ride-through \\
        HVDC  & High-voltage direct current \\
      \end{tabularx}}
    \end{column}
    \begin{column}{0.48\textwidth}
      {\small
      \begin{tabularx}{\linewidth}{@{}l@{\hspace{0.8em}}X@{}}
        IEEE  & Institute of Electrical and Electronics Engineers \\
        PSCAD & Power Systems Computer Aided Design \\
        PSLF  & Positive Sequence Load Flow \\
        PSS/E & Power System Simulator for Engineering \\
        RMS   & Root mean square \\
      \end{tabularx}}
    \end{column}
  \end{columns}
\end{frame}

% ---- Overview -----------------------------------------------
\begin{frame}{Design process\\System studies}
  \noindent{\color{atGray}\small Results of each study feed back into the steady-state design.}
  \vspace{0.6em}

  \centering
  \input{\ssfig{design-process}}
\end{frame}

% ---- 1 Steady state -----------------------------------------
\begin{frame}{Steady state\\Operating points, ratings and losses}
  \begin{columns}[T,onlytextwidth]
    \begin{column}{0.56\textwidth}
      \begin{block}{Calculation, no simulation tool}
        Algebraic calculation of the converter station at its steady-state operating points.
      \end{block}
      \begin{itemize}
        \setlength{\itemsep}{0.4em}
        \item Operating points across the power, AC voltage and DC voltage range
        \item Steady-state rating of the components
        \item Losses per component and for the station
        \item Starting point for all following studies
      \end{itemize}
    \end{column}
    \begin{column}{0.40\textwidth}
      \ssstudytable{MATLAB, Python}{no time domain}{Operating points, component ratings, losses}
    \end{column}
  \end{columns}
\end{frame}

% ---- 2 Transient --------------------------------------------
\begin{frame}{Transient simulation\\Converter internal faults}
  \begin{columns}[T,onlytextwidth]
    \begin{column}{0.56\textwidth}
      \begin{block}{EMT simulation, detailed model}
        Transient behaviour of the system during converter internal faults.
      \end{block}
      \begin{itemize}
        \setlength{\itemsep}{0.4em}
        \item Converter internal faults, e.g.\ valve and converter bus faults
        \item Design of the equipment for the transient stresses
        \item Design of protective equipment such as surge arresters
        \item Returns updated ratings to the steady-state design
      \end{itemize}
    \end{column}
    \begin{column}{0.40\textwidth}
      \ssstudytable{PSCAD, EMTP}{0.5--10\,µs}{Equipment stresses, arrester ratings}
    \end{column}
  \end{columns}
\end{frame}

% ---- 3 Dynamic ----------------------------------------------
\begin{frame}{Dynamic simulation\\AC faults and grid code compliance}
  \begin{columns}[T,onlytextwidth]
    \begin{column}{0.56\textwidth}
      \begin{block}{EMT simulation, control-oriented model}
        Dynamic behaviour of the system during AC faults.
      \end{block}
      \begin{itemize}
        \setlength{\itemsep}{0.4em}
        \item AC faults at and near the point of connection
        \item Grid code compliance defined by e.g.\ IEEE~2800
        \item FRT performance and control response
        \item Returns control and FRT limits to the design
      \end{itemize}
    \end{column}
    \begin{column}{0.40\textwidth}
      \ssstudytable{PSCAD, EMTP, Simulink}{10--20\,µs}{Control performance, FRT behaviour, compliance evidence}
    \end{column}
  \end{columns}
\end{frame}

% ---- 4 RMS --------------------------------------------------
\begin{frame}{RMS model\\Integration in the AC grid}
  \begin{columns}[T,onlytextwidth]
    \begin{column}{0.56\textwidth}
      \begin{block}{Phasor simulation, grid-level model}
        Simulation of the integration of the HVDC system in the AC grid.
      \end{block}
      \begin{itemize}
        \setlength{\itemsep}{0.4em}
        \item HVDC system as part of the wider AC network
        \item Model handed over to the grid operator for planning studies
        \item Returns grid integration requirements to the design
      \end{itemize}
    \end{column}
    \begin{column}{0.40\textwidth}
      \ssstudytable{PSS/E, PSLF, PowerFactory}{$\approx$\,1\,ms}{RMS model for grid studies}
    \end{column}
  \end{columns}
\end{frame}
